\section{Agent Roles}
\label{sec:agent_roles}

Prefiltering agents will invoke tools to evaluate the quality of input frames, filtering out frames with poor quality before passing them to the Actor agent.

The Actor Agent's job is to mutate the configuration state of the Gaussian Splatting reconstruction. The actor holds the current config file containing parameters such as number of Gaussians, densification interval, learning rates for position/opacity/scale/rotation, spherical harmonics degree, pruning thresholds, and camera pose refinement flags and applies edits to it based on structured feedback it receives from the Critic. Instead of doing generation from scratch for each iteration, the Actor applies edits such as "increase the densification interval from 100 to 150 steps. The Actor receives the Critic's structured feedback, reasons about which config parameters are causally responsible for the reported artifact, proposes a change to the config, applies it, triggers a re-train or re-optimize pass via gsplat, and then passes the newly rendered output back to the Critic for the next round. This gradient-free optimization loop over the config space, driven by visual and geometric feedback, is better than a loss function, because many failure modes in 3DGS (e.g. floaters, needle Gaussians, blurry textures, background collapse) are not easily encoded as a differentiable objective function.

The Critic Agent's job is to look at the rendered 3DGS output and produce structured, actionable feedback that the Actor can use to change specific parameters. The Critic will inspect the scene from multiple angles. It examines close-up views of problematic regions, depth maps to detect geometric inconsistencies, opacity visualizations to catch floaters, and normal maps to identify surface reconstruction failures

For the Critic agent, we propose using a powerful VLM such as GPT-4o or Claude. We take inspiration from frameworks like Feature4X, where gpt-4o act as high-level agents for executing and refining complex scene manipulations. In our architecture, the VLM Critic functions as an expert diagnostic layer. By analyzing a multi-view grid of renders alongside depth and opacity buffers, the agent can pinpoint specific optimization failures. For instance, the Critic might identify that “inward-pointing needle Gaussians on a ceiling” suggest a lack of upward-facing or overhead views.

The design constraint is that the Critic's output must be structured enough for the Actor to parse and act on a typed schema rather than a free-text paragraph the Actor has to interpret. This structured output lets us replay the full sequence of Critic observations and Actor edits to understand exactly how the scene evolved across iterations.